\documentclass[10pt,a4paper]{amsart}
\usepackage[margin=19mm]{geometry}
\usepackage{amsmath,amssymb,mathtools}
\usepackage[scr=rsfs,cal=euler]{mathalfa}
\usepackage{xfrac}
\usepackage[unicode,hidelinks,bookmarks=false]{hyperref}
\newcommand{\ie}{i.e.\ }
\newcommand{\cf}{cf.\ }
\newcommand{\resp}{respectively\ }
\newcommand{\Subsection}{\subsection}
\providecommand{\nobreakdash}{\nobreak\hbox{-}}
\setlength{\emergencystretch}{2em}
\multlinegap0pt
\usepackage[shortlabels]{enumitem}
\usepackage{xcolor}
\usepackage{amssymb}
\usepackage{algorithm}\usepackage{algpseudocode}
\usepackage{bm}
\usepackage{mathtools}
\newtheorem{theorem}{Theorem}
\def\lp{\left(}
\def\rp{\right)}
\def\sN{\mathcal{N}}
\def\sS{\mathcal{S}}
\def\sT{\mathcal{T}}
\title{\textbf{Localized Distributional Robustness in Submodular Multi-Task Subset Selection}}
\author{Ege C. Kaya, Abolfazl Hashemi}
\date{Source: arXiv:2404.03759v3 (CC BY 4.0)}
\begin{document}
\maketitle

\noindent\emph{Source and licence.} \textbf{Localized Distributional Robustness in Submodular Multi-Task Subset Selection}. Version arXiv:2404.03759v3 (CC BY 4.0), distributed by arXiv under the Creative Commons Attribution 4.0 International licence, http://creativecommons.org/licenses/by/4.0/, as recorded in the arXiv licence metadata for this exact version and shown on the abstract page https://arxiv.org/abs/2404.03759v3. Full licence notice: \texttt{licenses/arxiv-2404.03759-CC-BY-4.0.txt}.\par\smallskip

\noindent\textit{Editorial scope.} Counted unit: the article's theorem prop:main (source0.tex lines 547--553). The article's complete original proof of it is delivered with it (source0.tex lines 554--616, one \texttt{proof} environment of 63 lines). No mathematics is written, repaired or re-derived: the statement and the proof are the article's own lines, in the article's own notation, and no retained line is substituted or re-flowed.  No further source-local context is needed: every cross-reference of every retained line resolves inside the two delivered spans.  Nothing was omitted from any retained span, so the package carries no omission record.  No retained line contains a citation, so the wrapper carries no bibliography and no bibliography entry.  No retained line contains a figure environment, an image-inclusion command or drawing code, so no image is delivered and the package declares no asset dependency.  Editorial formatting only, in addition: an AMS \texttt{amsart} preamble that restates the article's own declarations and macros used by the retained lines, namely the environments \texttt{theorem} on separate counters, together with the standard packages those lines need.  The retained environments print in this document as Theorem 1 in order of appearance, whereas the article prints the same items with the numbering of its own full document.  The complete native source of the article as distributed in the arXiv e-print for this exact version is bundled unchanged as \texttt{sources/157238/source0.tex}.
\begin{theorem}\label{prop:main}
The set function $G(\sS)$ can be expressed as the composition of two functions, $G(S) = g(h(\sS))$, where the set function
\begin{equation}
h(\sS) \vcentcolon= \sum_{i=1}^n Q_i \lp1-\exp\lp-\frac{f_i(\sS)}{\lambda}\rp\rp
\end{equation}
is \textbf{(i)} normalized, \textbf{(ii)} monotone nondecreasing, and \textbf{(iii)} submodular, and the function $g(x) \vcentcolon= -\lambda\log(1-x)$ is \textbf{(iv)} monotone increasing, \textbf{(v)} convex, and \textbf{(vi)} Lipschitz continuous in $\operatorname{dom}g = [0, h(\sS^\ast)]$.
\end{theorem}

\begin{proof}
\begin{enumerate}[label=\bfseries(\roman*), leftmargin=*, itemsep=2pt]
\item We have
\begin{equation}
\begin{split}
h(\emptyset) = \sum_{i=1}^nQ_i\lp1-\exp{\lp-\frac{f_i(\emptyset)}{\lambda}\rp}\rp 
= \sum_{i=1}^nQ_i\lp1-\exp{\lp0\rp}\rp = 0.
\end{split}
\end{equation}
\item Let $\sS \subseteq \mathcal{T} \subseteq \sN$. Then, for each $i \in [n],$ we have $f_i(\sS) \leq f_i(\mathcal{T}).$ Hence, $-f_i(\sS)/\lambda \geq - f_i(\mathcal{T})/\lambda,$ and $\exp(-f_i(\sS)/\lambda) \geq \exp(-f_i(\mathcal{T})/\lambda),$ thanks to the monotone increasing property of the $\exp$ function and the positivity of $\lambda$. This means that 
\begin{equation}
1-\exp\lp -\frac{f_i(\sS)}{\lambda}\rp \leq  1-\exp\lp -\frac{f_i(\mathcal{T})}{\lambda}\rp,
\end{equation}
and because $Q_i \geq 0$ for all $i \in [n],$ 
\begin{equation}
\begin{split}
\sum_{i=1}^nQ_i\lp1-\exp\lp -\frac{f_i(\sS)}{\lambda}\rp \rp \leq \sum_{i=1}^nQ_i\lp1-\exp\lp -\frac{f_i(\mathcal{T})}{\lambda}\rp\rp,
\end{split}
\end{equation}
that is, $h(\sS) \leq h(\mathcal{T})$.
\item Let $\sS \subseteq \mathcal{T} \subset \sN$, $e \in \sN \setminus \mathcal{T}$. For all $i \in [n]$, we have
\begin{equation}
\dfrac{f_i(\sS \cup \{e\}) - f_i(\sS)}{\lambda} \geq \dfrac{f_i(\mathcal{T} \cup \{e\}) - f_i(\mathcal{T})}{\lambda}.
\end{equation}
Now, because $x\mapsto 1-\exp(-x)$ is a monotone increasing function, we have
\begin{equation}\label{eq:multiplythis}
\begin{split}
1-\exp\lp\dfrac{-f_i(\sS \cup \{e\}) + f_i(\sS)}{\lambda}\rp \geq 1-\exp\lp\dfrac{-f_i(\mathcal{T} \cup \{e\}) + f_i(\mathcal{T})}{\lambda} \rp.
\end{split}
\end{equation}
Furthermore, because each $f_i$ is monotone nondecreasing, $f_i(\sS) \leq f_i(\mathcal{T})$, and since $x\mapsto \exp(-x/\lambda)$ is a monotone decreasing function for $\lambda > 0$, we have
\begin{equation}
\exp\lp \frac{-f_i(\sS)}{\lambda} \rp \geq \exp\lp \frac{-f_i(\mathcal{T})}{\lambda} \rp.
\end{equation}
Both of these quantities are nonnegative. Then, multiplying the left-hand side of \eqref{eq:multiplythis} by $\exp(-f_i(\sS)/\lambda)$ and the right-hand side by $\exp(-f_i(\mathcal{T})/\lambda)$ preserves the inequality, and we obtain
\begin{equation}
\begin{split}
\exp\lp \frac{-f_i(\sS)}{\lambda} \rp - \exp\lp \frac{-f_i(\sS \cup \{e\})}{\lambda} \rp  \geq \exp\lp \frac{-f_i(\mathcal{T})}{\lambda} \rp - \exp\lp \frac{-f_i(\sT \cup \{e\})}{\lambda} \rp.
\end{split}
\end{equation}
This demonstrates that the set function $\sS \mapsto -\exp (-f_i(\sS)/\lambda)$ is submodular for all $i \in [n]$. The addition of the constant $1$, multiplication with the nonnegative $Q_i$ and the summation through the indices $i \in [n]$ preserve submodularity, hence,
\begin{equation}
h(\sS) = \sum_{i=1}^n Q_i \lp 1 - \exp\lp \dfrac{-f_i (\sS)}{\lambda}\rp\rp
\end{equation}
is submodular.
\item It suffices to note that the first-order derivative of $g$,
\begin{equation}
\dfrac{\mathrm{d}g}{\mathrm{d}x} = \dfrac{\lambda}{1-x} > 0,
\end{equation}
for all $\lambda >0$ and $x \in [0, h(\sS^\ast)]$. 
% Furthermore $\frac{\mathrm{d}g}{\mathrm{d}x}$ attains its maximum value at $x=h(S^\ast)$.
\item The second-order derivative of g,
\begin{equation}
\dfrac{\mathrm{d}^2g}{\mathrm{d}x^2} = \dfrac{\lambda}{(1-x)^2} > 0,
\end{equation}
for all $\lambda >0$ and $x \in [0, h(\sS^\ast)]$, hence $g$ is a convex function.
\item As per \textbf{(v)}, the first-order derivative of $g$ is a monotone increasing function over its domain, attaining its maximum value at $x = h(\sS^\ast)$. Hence, 
\begin{equation}
\left\lvert\dfrac{\mathrm{d}g}{\mathrm{d}x}\right\rvert \leq \dfrac{\lambda}{1-h(\sS^\ast)},
\end{equation}
making $g$ Lipschitz continuous with constant $\lambda/(1-h(\sS^\ast))$.
\end{enumerate}
\end{proof}

\end{document}
